\chapter{完全信息静态博弈}\label{chap:01}

% ---------------------------------------------------------------------------
\exer{1.1}
什么是博弈的标准式？在博弈的标准式中，什么是严格劣战略？什么是一个纯战略纳什均衡？

\sol
\textbf{（1）博弈的标准式（normal form）}：一个 $n$ 人博弈的标准式表述由以下三个要素组成：
\begin{enumerate}[leftmargin=2em]
  \item 参与人集合：$\{1,2,\dots,n\}$；
  \item 每个参与人的战略空间 $S_i$（即其可选择的全部战略的集合）；
  \item 每个参与人的收益函数 $u_i:S_1\times S_2\times\cdots\times S_n\to\mathbb{R}$，它把每一个战略组合 $s=(s_1,\dots,s_n)$ 映射为参与人 $i$ 的收益 $u_i(s)$。
\end{enumerate}
通常记为 $G=\{S_1,\dots,S_n;u_1,\dots,u_n\}$。

\textbf{（2）严格劣战略（strictly dominated strategy）}：在给定的标准式博弈中，若存在另一个战略 $s_i'\in S_i$（可以是纯战略，也可以是混合战略），使得对对手所有可能的战略组合 $s_{-i}\in S_{-i}$，都有
\[
  u_i(s_i',s_{-i})>u_i(s_i,s_{-i}),
\]\indent
则称 $s_i$ 为参与人 $i$ 的严格劣战略。理性参与人不会选择严格劣战略。

\textbf{（3）纯战略纳什均衡}：战略组合 $s^*=(s_1^*,\dots,s_n^*)$ 是一个纯战略纳什均衡，当且仅当对每个参与人 $i$，在其他参与人都选择 $s_{-i}^*$ 的情况下，$s_i^*$ 是 $i$ 的最优反应，即
\[
  u_i(s_i^*,s_{-i}^*)\ge u_i(s_i,s_{-i}^*),\qquad \forall s_i\in S_i.
\]

% ---------------------------------------------------------------------------
\exer{1.2}
在以下博弈的标准式中，哪些战略不会被重复剔除严格劣战略所剔除？纯战略纳什均衡又是什么？

\begin{paytab}
\begin{tabular}{ccccc}
\multicolumn{2}{c}{} & \multicolumn{3}{c}{\pthead{参与人 2}}\\
\cline{3-5}
\multicolumn{2}{c}{} & $L$ & $C$ & $R$\\[1pt]
\hline
\multirow{3}{*}{参与人 1} & $T$ & $2,0$ & $1,1$ & $4,2$\\[3pt]
\cline{2-5}
 & $M$ & $3,4$ & $1,2$ & $2,3$\\[3pt]
\cline{2-5}
 & $B$ & $1,3$ & $0,2$ & $3,0$\\[3pt]
\hline
\end{tabular}
\end{paytab}

\sol
先求各参与人的最优反应：
\begin{enumerate}[leftmargin=2em]
  \item 对参与人 1（行）：若参与人 2 选 $L$，最优为 $M$（得 3）；选 $C$，$T$、$M$ 并列最优（得 1）；选 $R$，最优为 $T$（得 4）。
  \item 对参与人 2（列）：若参与人 1 选 $T$，最优为 $R$（得 2）；选 $M$，最优为 $L$（得 4）；选 $B$，最优为 $L$（得 3）。
\end{enumerate}

\textbf{第一步：} 检查参与人 1 的 $B$ 是否被某个混合战略严格劣于。取以概率 $3/4$ 选 $T$、概率 $1/4$ 选 $M$ 的混合战略，则
当参与人 2 选择 $L$ 时，混合战略的收益为
\[
  \frac{3}{4}\cdot 2+\frac{1}{4}\cdot 3
  =\frac{9}{4}>1.
\]
\indent 当参与人 2 选择 $C$ 时，混合战略的收益为
\[
  \frac{3}{4}\cdot 1+\frac{1}{4}\cdot 1=1>0.
\]
\indent 当参与人 2 选择 $R$ 时，混合战略的收益为
\[
  \frac{3}{4}\cdot 4+\frac{1}{4}\cdot 2
  =\frac{7}{2}>3.
\]
\indent 故 $B$ 是严格劣战略，可剔除。

\textbf{第二步：} 剔除 $B$ 后，剩下 $2\times3$ 博弈。对参与人 2，$C$ 被某个 $L$、$R$ 的混合战略严格劣于。取以概率 $1/4$ 选 $L$、概率 $3/4$ 选 $R$，则
当参与人 1 选择 $T$ 时，混合战略的期望收益为
\[
  \frac{1}{4}\cdot 0+\frac{3}{4}\cdot 2
  =\frac{3}{2}>1.
\]
\indent 当参与人 1 选择 $M$ 时，混合战略的期望收益为
\[
  \frac{1}{4}\cdot 4+\frac{3}{4}\cdot 3
  =\frac{13}{4}>2.
\]\indent
故 $C$ 可剔除。

\textbf{第三步：} 剩余 $2\times2$ 博弈中，$T$、$M$、$L$、$R$ 均不被严格劣于，因此
\[
  \boxed{\text{重复剔除严格劣战略后 surviving 的战略为：}\{T,M\}\times\{L,R\}}.
\]

\Needspace{0.30\textheight}
在剔除 $B$、$C$ 后的 $2\times2$ 支付矩阵中，用划线标出各参与人的最佳反应。每个格子的收益次序为“参与人 1，参与人 2”：
\begin{paytab}
\normalsize
\renewcommand{\arraystretch}{1.35}
\setlength{\tabcolsep}{8pt}
\begin{tabular}{ccccc}
\multicolumn{2}{c}{} & \multicolumn{2}{c}{\pthead{参与人 2}}\\
\cline{3-4}
\multicolumn{2}{c}{} & $L$ & $R$\\[1pt]
\hline
\multirow{2}{*}{参与人 1} & $T$ & $2,\ 0$ & $\underline{4},\ \underline{2}$\\[3pt]
\cline{2-4}
 & $M$ & $\underline{3},\ \underline{4}$ & $2,\ 3$\\[3pt]
\hline
\end{tabular}
\end{paytab}
第一个收益被划线表示参与人 1 对该列行动的最佳反应；第二个收益被划线表示参与人 2 对该行行动的最佳反应。双方收益都被划线的格子对应纳什均衡，因此
\[
  \boxed{(M,L)\text{ 和 }(T,R)}.
\]

% ---------------------------------------------------------------------------
\exer{1.3}
两个人就如何分配一元钱进行谈判，双方同时提出各自希望得到的份额，分别为 $s_1$ 和 $s_2$，且 $0\le s_1,s_2\le1$。若 $s_1+s_2\le1$，则二人分别得到他们所要的一份；如果 $s_1+s_2>1$，则两个人均一无所获。求出此博弈的纯战略纳什均衡。

\sol
设参与人 $i$ 提出 $s_i\in[0,1]$，其收益为
\[
  u_i(s_i,s_j)=\begin{cases}
    s_i, & s_i+s_j\le1,\\
    0, & s_i+s_j>1.
  \end{cases}
\]

\textbf{情形 1：}$s_1+s_2=1$。此时 $u_1=s_1$，$u_2=s_2$。对参与人 1，若提高份额至 $s_1'>s_1$，则 $s_1'+s_2>1$，收益降为 0；若降低至 $s_1'<s_1$，收益变为 $s_1'<s_1$。故 $s_1$ 是最优反应。同理 $s_2$ 也是最优反应。因此所有满足 $s_1+s_2=1$ 的战略组合都是纯战略 NE。

\textbf{情形 2：}$s_1+s_2<1$。此时两人收益均为正，且均可单方面提高自己的报价而不触发冲突，从而严格增加收益。故不存在 NE。

\textbf{情形 3：}$s_1+s_2>1$ 但 $(s_1,s_2)\ne(1,1)$。此时至少有一人收益为 0，且该参与人可通过把报价降到 $1-s_j>0$ 而获得正收益。故不是 NE。

\textbf{情形 4：}$(s_1,s_2)=(1,1)$。两人收益均为 0。参与人 1 的任何偏离 $s_1'<1$ 都会使 $s_1'+s_2=s_1'+1\le1$ 仅当 $s_1'=0$，收益仍为 0；任何 $s_1'>1$ 不可行。故 $(1,1)$ 也是 NE。

综上，纯战略纳什均衡集合为
\[
  \boxed{\{(s_1,s_2):s_1+s_2=1,\ 0\le s_1,s_2\le1\}\ \cup\ \{(1,1)\}}.
\]

% ---------------------------------------------------------------------------
\exer{1.4}
假定古诺的寡头垄断模型中有 $n$ 个企业，令 $q_i$ 代表企业 $i$ 的产量，且 $Q=q_1+\cdots+q_n$ 表示市场总产量，$p$ 表示市场出清价格，并假设反需求函数由 $p(Q)=a-Q$（设 $Q<a$，其他情况下 $p=0$）给出。并设企业 $i$ 生产出 $q_i$ 的总成本 $C_i(q_i)=cq_i$，即没有固定成本，且边际成本为常数 $c$，这里我们设 $c<a$。根据古诺的假定，企业同时就产量进行决策。求出博弈的纳什均衡。当 $n$ 趋于无穷时，将会发生什么情况？

\sol
企业 $i$ 的利润函数为
\[
  \pi_i(q_i,q_{-i})=\bigl[p(Q)-c\bigr]q_i=(a-q_i-q_{-i}-c)q_i,
\]\indent
其中 $q_{-i}=\sum_{j\ne i}q_j$。

\textbf{（1）求纳什均衡。} 对 $q_i$ 求一阶条件：
\[
  \frac{\partial\pi_i}{\partial q_i}=a-2q_i-q_{-i}-c=0,
\]\indent
得到企业 $i$ 的反应函数
\[
  q_i=\frac{a-q_{-i}-c}{2}.
\]\indent
由于所有企业同质，均衡时 $q_1^*=q_2^*=\cdots=q_n^*=q^*$，于是 $q_{-i}^*=(n-1)q^*$。代入反应函数：
\[
  q^*=\frac{a-(n-1)q^*-c}{2}
  \quad\Longrightarrow\quad
  q^*=\frac{a-c}{n+1}.
\]\indent
因此纳什均衡为
\[
  \boxed{(q_1^*,\dots,q_n^*)=\Bigl(\frac{a-c}{n+1},\dots,\frac{a-c}{n+1}\Bigr)}.
\]\indent
市场总产量与均衡价格为
\[
  Q^*=\frac{n(a-c)}{n+1},\qquad p^*=a-Q^*=\frac{a+nc}{n+1}.
\]\indent
二阶条件 $-2<0$ 保证该解为利润最大值点。

\textbf{（2）$n\to\infty$ 的极限。}
\[
  \lim_{n\to\infty}Q^*=a-c,\qquad \lim_{n\to\infty}p^*=c.
\]\indent
即价格收敛于边际成本 $c$，市场趋于完全竞争市场。

% ---------------------------------------------------------------------------
\exer{1.5}
考虑以下两个古诺双头垄断模型的战略空间有限的情况。

第一，假定每个企业必须选择要么生产垄断产出的一半 $q_m/2=(a-c)/4$，要么生产古诺均衡产量 $q_c=(a-c)/3$，任何其他产量都是不允许的。证明这一非此即彼的博弈是一个囚徒困境式的问题：每一个企业都有一个严格劣战略，并且在均衡状态下，每一个企业的福利都要比他们相互合作时下降。

第二，假设每个企业可以选择 $q_m/2$ 或 $q_c$，或第三种产量 $q'$，求出一个 $q'$ 的值，使得这一博弈在以下方面等价于第 1.2.A 节中的古诺模型，即 $(q_c,q_c)$ 是唯一的纳什均衡，并且在均衡状态下，每一企业的福利都比他们相互合作时要低，但两个企业都没有严格劣战略。

\sol
记 $d=a-c>0$。两种可选产量分别为
\[
  q_m/2=\frac{d}{4},\qquad q_c=\frac{d}{3}.
\]

\textbf{（1）二战略博弈。} 企业 $i$ 的利润为 $\pi_i=(a-q_i-q_j-c)q_i=dq_i-(q_i+q_j)q_i$。计算四种情形下的利润：
\begin{enumerate}[leftmargin=2em]
  \item 双方都选 $q_m/2$：$Q=d/2$，$p=c+d/2$，$\pi_i=\dfrac{d}{4}\cdot\dfrac{d}{2}=\dfrac{d^2}{8}$。
  \item 企业 $i$ 选 $q_m/2$，企业 $j$ 选 $q_c$：$Q=7d/12$，$p=c+5d/12$，
    \[
      \pi_i=\frac{d}{4}\cdot\frac{5d}{12}=\frac{5d^2}{48},\qquad
      \pi_j=\frac{d}{3}\cdot\frac{5d}{12}=\frac{5d^2}{36}.
    \]\indent
  \item 双方都选 $q_c$：$Q=2d/3$，$p=c+d/3$，$\pi_i=\dfrac{d}{3}\cdot\dfrac{d}{3}=\dfrac{d^2}{9}$。
\end{enumerate}
支付矩阵如下：

\begin{paytab}
\begin{tabular}{cccc}
\multicolumn{2}{c}{} & \multicolumn{2}{c}{\pthead{企业 2}}\\
\cline{3-4}
\multicolumn{2}{c}{} & $q_m/2$ & $q_c$\\[2pt]
\hline
\paybox{1.2cm}{企业 1} & \paybox{1.5cm}{$q_m/2$} & \paybox{2.1cm}{$\dfrac{d^2}{8},\ \dfrac{d^2}{8}$} & \paybox{2.1cm}{$\dfrac{5d^2}{48},\ \dfrac{5d^2}{36}$}\\
\hline
\paybox{1.2cm}{企业 1} & \paybox{1.5cm}{$q_c$} & \paybox{2.1cm}{$\dfrac{5d^2}{36},\ \dfrac{5d^2}{48}$} & \paybox{2.1cm}{$\dfrac{d^2}{9},\ \dfrac{d^2}{9}$}\\
\hline
\end{tabular}
\end{paytab}

逐项比较两种产量对应的利润：
\begin{align*}
  \frac{d^2}{8}&=\frac{6d^2}{48}<\frac{5d^2}{36},\\
  \frac{5d^2}{48}&<\frac{d^2}{9}.
\end{align*}\indent
因此，对企业 1 和企业 2 而言，无论对方选择哪种产量，选择 $q_c$ 的利润都更高。故 $q_m/2$ 是严格劣战略。

剔除 $q_m/2$ 后，唯一剩余的战略组合为 $(q_c,q_c)$，这就是唯一的纳什均衡。

合作结果 $(q_m/2,q_m/2)$ 与均衡结果 $(q_c,q_c)$ 下，每个企业的利润分别为
\[
  \pi_i^{C}=\frac{d^2}{8},\qquad
  \pi_i^{NE}=\frac{d^2}{9}<\frac{d^2}{8}.
\]
\indent 因此，均衡时福利低于合作水平——这正是囚徒困境的结构。

\textbf{（2）加入第三种产量后，题目所列条件是否可同时满足？} 令 $y=q'/d$。对任意自身产量 $x$、对手产量 $z$（均以 $d=a-c$ 为单位），利润为
\[
  \pi(x,z)=d^2x(1-x-z).
\]\indent
比较 $q_c=d/3$ 与 $q_m/2=d/4$：
\[
  \pi\!\left(\frac13,z\right)-\pi\!\left(\frac14,z\right)
  =\frac{d^2}{144}(5-12z).
\]

若 $y<5/12$，则对手可选的三个产量 $1/4,1/3,y$ 全都小于 $5/12$。上式在每一种对手选择下都严格为正，故 $q_c$ 严格支配 $q_m/2$，与“没有严格劣战略”矛盾。

若 $y>5/12$，则对任意对手产量 $z\in\{1/4,1/3,y\}$，
\[
  \pi\!\left(\frac13,z\right)-\pi(y,z)
  =d^2\left(\frac13-y\right)\left(\frac23-y-z\right)>0,
\]\indent
因为两因子均为负数。此时 $q_c$ 严格支配 $q'$，仍与题设要求矛盾。

边界 $y=5/12$ 也不能满足“唯一纳什均衡”：对手选 $q_m/2$ 时，$q_c$ 与 $q'$ 都是最佳反应；对手选 $q'$ 时，$q_m/2$ 与 $q_c$ 都是最佳反应。因此 $(q',q_m/2)$ 和 $(q_m/2,q')$ 也是纯战略纳什均衡。

所以按当前题面，\textbf{不存在}一个 $q'$ 能同时满足“$(q_c,q_c)$ 是唯一纳什均衡”与“双方均无严格劣战略”。例如 $q'=d/2$ 确实使 $(q_c,q_c)$ 成为唯一均衡，但此时 $q'$ 被 $q_c$ 严格支配，仍不满足题目最后一个条件。原题这一组要求彼此不相容，可能存在印刷或翻译问题；在未核实勘误前，不应把某个 $q'$ 写成完全满足全部条件的答案。

% ---------------------------------------------------------------------------
\exer{1.6}
考虑在古诺双头垄断模型中，反需求函数为 $p(Q)=a-Q$，但两企业有不同的边际成本，企业 1 为 $c_1$，企业 2 为 $c_2$，求出当每个企业 $0<c_i<a/2$ 时的纳什均衡。如果 $c_1<c_2<a$ 但 $2c_2>a+c_1$，纳什均衡又有什么变化？

\sol
企业 $i$ 的利润函数为
\[
  \pi_i=(a-q_i-q_j-c_i)q_i,\qquad i=1,2,\ j\ne i.
\]\indent
一阶条件给出反应函数
\[
  q_i=R_i(q_j)=\frac{a-q_j-c_i}{2},\qquad i=1,2.
\]

\textbf{（1）$0<c_i<a/2$ 时的内部解。} 联立两反应函数：
\begin{align*}
  q_1&=\frac{a-q_2-c_1}{2},\\
  q_2&=\frac{a-q_1-c_2}{2}.
\end{align*}\indent
解得
\[
  \boxed{q_1^*=\frac{a-2c_1+c_2}{3},\qquad q_2^*=\frac{a-2c_2+c_1}{3}}.
\]\indent
由 $c_i<a/2$ 易知 $q_i^*>0$，且二阶条件 $-2<0$ 保证是最优。故内部纳什均衡为 $(q_1^*,q_2^*)$。

\textbf{（2）$c_1<c_2<a$ 且 $2c_2>a+c_1$ 时的边界解。} 此时 $c_2$ 足够高，使得
\[
  a-2c_2+c_1<0,
\]\indent
即企业 2 在标准反应函数下的“意愿产量”为负。由于产量非负，企业 2 的最优反应为 $q_2=0$，只要企业 1 的产量足够大。企业 1 面对 $q_2=0$ 时选择垄断产量
\[
  q_1^*=\frac{a-c_1}{2}.
\]\indent
需要验证企业 2 确实不愿偏离 $q_2=0$：当企业 1 生产 $q_1^*=(a-c_1)/2$ 时，企业 2 生产任意正产量 $q_2>0$ 的利润为
\[
  \pi_2=q_2\Bigl(a-\frac{a-c_1}{2}-q_2-c_2\Bigr)
       =q_2\Bigl(\frac{a+c_1}{2}-c_2-q_2\Bigr).
\]\indent
由 $2c_2>a+c_1$ 知括号内首项为负，故对任意 $q_2>0$ 有 $\pi_2<0$，而 $q_2=0$ 时利润为 0。因此企业 2 没有动力偏离。

综上，当 $2c_2>a+c_1$ 时，纳什均衡为企业 2 退出市场，企业 1 独占垄断产量：
\[
  \boxed{(q_1^*,q_2^*)=\Bigl(\frac{a-c_1}{2},\ 0\Bigr)}.
\]

% ---------------------------------------------------------------------------
\exer{1.7}
在第 1.2.B 中，我们分析了产品有差异的贝特兰德双头垄断模型。同质产品的情况下结论是十分明显的。假设 $p_i<p_j$ 时，消费者对企业 $i$ 产品的需求为 $a-p_i$；$p_i>p_j$ 时为 $0$；$p_i=p_j$ 时为 $(a-p_i)/2$。同时假设不存在固定成本，且边际成本为常数 $c$，这里 $c<a$。证明如果企业同时选择价格，则唯一的纳什均衡就是每个企业的定价均为 $c$。

\sol
设企业 $i$ 定价 $p_i\ge c$（定价低于边际成本会使每单位亏损，故理性企业不会这么做）。企业 $i$ 的利润函数为
\[
  \pi_i(p_i,p_j)=\begin{cases}
    (p_i-c)(a-p_i), & p_i<p_j,\\[6pt]
    \dfrac{(p_i-c)(a-p_i)}{2}, & p_i=p_j,\\[10pt]
    0, & p_i>p_j.
\end{cases}
\]

\textbf{第一步：} 证明 $p_i^*=p_j^*=c$ 是 NE。当对方定价为 $c$ 时，若企业 $i$ 也定价 $c$，利润为 0；若提价至 $p_i>c$，则 $p_i>p_j=c$，需求为 0，利润仍为 0；若降价至 $p_i<c$，则每单位亏损，利润为负。故 $p_i=c$ 是最优反应。由对称性，$(c,c)$ 是 NE。

\textbf{第二步：} 证明唯一性。设 $(p_1^*,p_2^*)$ 为 NE。若 $p_1^*>c$，则企业 2 只要定价 $p_2\in(c,p_1^*)$ 即可获得全部市场需求和正利润，因此 $p_2^*\ne p_1^*$ 且 $p_2^*<p_1^*$。但此时企业 1 可略微降价至 $p_1^*-\varepsilon>p_2^*$？不对，若 $p_2^*<p_1^*$，企业 1 可降价到 $p_2^*-\varepsilon$ 以下重新赢得市场，获得正利润。因此 $p_1^*>c$ 不能是均衡。同理 $p_2^*>c$ 也不能。若 $p_1^*=p_2^*>c$，则每家企业各得一半市场；任一方略微降价即可夺取全部市场，利润几乎翻倍，故会偏离。因此唯一可能的 NE 是 $p_1^*=p_2^*=c$。

\[
  \boxed{(p_1^*,p_2^*)=(c,c)}.
\]

% ---------------------------------------------------------------------------
\exer{1.8}
设有一批选民在一个单位区间从左（$x=0$）至右（$x=1$）均匀分布，为一个职位参加竞选的每个候选人同时选择其竞选基地（即在 $x=0$ 到 $x=1$ 中间的一个点）。选民观察候选人的选择，然后每一投票人把票投给其基地离自己最近的候选人。比如，如果有两个候选人，他们分别在 $x_1=0.3$ 和 $x_2=0.6$ 选择基地，则处于 $x=0.45$ 左边的所有选民都会把票投给候选人 1，右边的人都会把票投给候选人 2，这样候选人 2 就可以得到 $55\%$ 的选票赢得这场选举。假设候选人只关心他能否当选——他们根本上一点都不关心其基地！如果有两个候选人，博弈的纯战略纳什均衡是什么？如果有三个候选人，求出一个纳什均衡。（假设选择同一个基地的候选人将平分这一基地可得的选票，得票最高的候选人不止一人时，谁当选由掷硬币来决定。）参见霍特林（Hotelling，1929）关于此类博弈的早期模型。

\sol
候选人只关心当选概率，因此一个位置是最佳反应当且仅当它没有给该候选人带来严格更高的获胜概率。

\textbf{（1）两个候选人。} 设两候选人分别位于 $x_1\le x_2$。
\begin{enumerate}[leftmargin=2em]
  \item 若 $x_1<x_2$，则候选人 1 得票 $(x_1+x_2)/2$，候选人 2 得票 $1-(x_1+x_2)/2$。若 $x_1+x_2<1$，候选人 1 得票不足一半，可向右微移直至超过中心而获胜；若 $x_1+x_2>1$，候选人 2 得票不足一半，可向左微移获胜。若 $x_1+x_2=1$ 但 $x_1\ne1/2$，两人平分选票，但任一方移向中心 $1/2$ 即可赢得多数，故会偏离。
  \item 若 $x_1=x_2\ne1/2$，两人平分选票，各以 $1/2$ 概率当选；任一方移向中心 $1/2$ 即可赢得多数，故会偏离。
  \item 若 $x_1=x_2=1/2$，两人仍各得一半选票，各以 $1/2$ 概率当选。任何一方单独偏离到 $x\ne1/2$ 都会使自己的得票低于 $1/2$ 而落败（例如左移一点只得到略小于 $1/2$ 的左半区）。因此没有有利可图的偏离。
\end{enumerate}
故两个候选人时，唯一的纯战略纳什均衡为
\[
  \boxed{x_1=x_2=\frac12}.
\]

\textbf{（2）三个候选人。} 下面给出一个纯战略纳什均衡，并验证它。

考虑位置组合
\[
  \boxed{(x_1,x_2,x_3)=\Bigl(\frac14,\frac34,\frac34\Bigr)}.
\]\indent
在该组合下，候选人 1（位于 $1/4$）获得从 $0$ 到 $(1/4+3/4)/2=1/2$ 的选票，即得票 $1/2$；候选人 2、3 均位于 $3/4$，他们平分从 $1/2$ 到 $1$ 的选票，每人得票 $1/4$。因此候选人 1 当选，候选人 2、3 落选。

\textbf{验证无偏离动机：}
\begin{enumerate}[leftmargin=2em]
  \item 候选人 1 已以概率 1 当选，这是最高可能当选概率，因此任何偏离都不会严格提高其当选概率（某些偏离可能使其仍当选，某些会使其落败，但不存在严格更优的位置）。
  \item 对候选人 2（或 3），若其偏离到 $z<3/4$：
    \begin{enumerate}
      \item 若 $z\in(1/4,3/4)$，则该候选人得票为 $(3/4-1/4)/2=1/4$，仍小于 $1/2$，无法获胜；
      \item 若 $z\le1/4$，其得票不超过 $1/4$，无法获胜。
    \end{enumerate}
    若其偏离到 $z>3/4$，则该候选人得票为 $1-(3/4+z)/2$。要使其超过 $1/2$，需 $z<1/4$，与 $z>3/4$ 矛盾，故无法获胜。
\end{enumerate}

因此 $(1/4,3/4,3/4)$ 是一个纯战略纳什均衡。由对称性，$(3/4,1/4,1/4)$ 也是一个 NE。

\cmt{三候选人时存在许多非对称 NE，例如一般地 $(x,y,y)$（$x<y$）满足 $x+y\ge2/3$ 且 $x+3y\ge2$ 时均为 NE；上述 $(1/4,3/4,3/4)$ 是最简洁的例子之一。}

% ---------------------------------------------------------------------------
\exer{1.9}
什么是标准式博弈的混合战略？什么是标准式博弈的混合战略纳什均衡？

\sol
\textbf{（1）混合战略（mixed strategy）}：在标准式博弈 $G=\{S_1,\dots,S_n;u_1,\dots,u_n\}$ 中，参与人 $i$ 的混合战略 $\sigma_i$ 是其纯战略空间 $S_i$ 上的一个概率分布。若 $S_i=\{s_{i1},\dots,s_{ik}\}$，则 $\sigma_i=(p_{i1},\dots,p_{ik})$，其中 $p_{ij}\ge0$ 且 $\sum_j p_{ij}=1$。纯战略是混合战略的特例（把全部概率集中在某一纯战略上）。

\textbf{（2）混合战略纳什均衡}：混合战略组合 $\sigma^*=(\sigma_1^*,\dots,\sigma_n^*)$ 是一个混合战略纳什均衡，当且仅当对每个参与人 $i$，
\[
  u_i(\sigma_i^*,\sigma_{-i}^*)\ge u_i(\sigma_i,\sigma_{-i}^*),\qquad \forall \sigma_i\in\Delta(S_i),
\]\indent
其中 $u_i(\sigma_i,\sigma_{-i}^*)$ 是参与人 $i$ 在混合战略下的期望收益。等价地，$\sigma_i^*$ 中的每一个以正概率出现的纯战略都必须是对手战略组合 $\sigma_{-i}^*$ 的最佳反应。

% ---------------------------------------------------------------------------
\exer{1.10}
证明在 1.1 节中所分析的 3 个标准式博弈——囚徒困境、图 1.1.1 和图 1.1.4 中，不存在混合战略纳什均衡。

\sol
下面分别证明三个博弈均不存在非退化的混合战略纳什均衡。

\textbf{（1）囚徒困境。} 支付矩阵为
\begin{paytab}
\begin{tabular}{cccc}
\multicolumn{2}{c}{} & \multicolumn{2}{c}{\pthead{囚徒 2}}\\
\cline{3-4}
\multicolumn{2}{c}{} & 沉默 & 坦白\\[2pt]
\hline
\multirow{2}{*}{囚徒 1} & 沉默 & $-1,-1$ & $-9,\phantom{-}0$\\[4pt]
\cline{2-4}
 & 坦白 & $\phantom{-}0,-9$ & $-6,-6$\\[4pt]
\hline
\end{tabular}
\end{paytab}

对两个囚徒而言，坦白都严格优于沉默：无论对方选什么，坦白的刑期都更短（收益更高）。因此任何包含沉默正概率的混合战略都被“永远坦白”严格占优，不可能在 NE 中出现。唯一可能的战略组合是（坦白，坦白），它是一个纯战略 NE，不是混合战略 NE。故不存在混合战略 NE。

\textbf{（2）图 1.1.1。} 该博弈为
\begin{paytab}
\begin{tabular}{ccccc}
\multicolumn{2}{c}{} & \multicolumn{3}{c}{\pthead{参与人 2}}\\
\cline{3-5}
\multicolumn{2}{c}{} & $L$ & $M$ & $R$\\[2pt]
\hline
\multirow{2}{*}{参与人 1} & $U$ & $1,0$ & $1,2$ & $0,1$\\[4pt]
\cline{2-5}
 & $D$ & $0,3$ & $0,1$ & $2,0$\\[4pt]
\hline
\end{tabular}
\end{paytab}

对参与人 2，$R$ 被 $M$ 严格占优（因为 $2>1$ 且 $1>0$），所以任何 NE 中 $R$ 的概率必为 0。剔除 $R$ 后得到图 1.1.2：
\begin{paytab}
\begin{tabular}{cccc}
\multicolumn{2}{c}{} & \multicolumn{2}{c}{\pthead{参与人 2}}\\
\cline{3-4}
\multicolumn{2}{c}{} & $L$ & $M$\\[2pt]
\hline
\multirow{2}{*}{参与人 1} & $U$ & $1,0$ & $1,2$\\[4pt]
\cline{2-4}
 & $D$ & $0,3$ & $0,1$\\[4pt]
\hline
\end{tabular}
\end{paytab}

此时对参与人 1，$D$ 被 $U$ 严格占优（$1>0$ 且 $1>0$），故 $D$ 的概率为 0。再剔除 $D$ 后得到图 1.1.3：
\begin{paytab}
\begin{tabular}{ccc}
 & \multicolumn{2}{c}{\pthead{参与人 2}}\\
\cline{2-3}
 & $L$ & $M$\\[2pt]
\hline
$U$ & $1,0$ & $1,2$\\[4pt]
\hline
\end{tabular}
\end{paytab}

最后对参与人 2，$L$ 被 $M$ 严格占优（$2>0$）。于是唯一 NE 为 $(U,M)$，是纯战略。因此图 1.1.1 不存在混合战略 NE。

\textbf{（3）图 1.1.4。} 该博弈为
\begin{paytab}
\begin{tabular}{ccccc}
\multicolumn{2}{c}{} & \multicolumn{3}{c}{\pthead{参与人 2}}\\
\cline{3-5}
\multicolumn{2}{c}{} & $L$ & $C$ & $R$\\[2pt]
\hline
\multirow{3}{*}{参与人 1} & $T$ & $0,4$ & $4,0$ & $5,3$\\[4pt]
\cline{2-5}
 & $M$ & $4,0$ & $0,4$ & $5,3$\\[4pt]
\cline{2-5}
 & $B$ & $3,5$ & $3,5$ & $6,6$\\[4pt]
\hline
\end{tabular}
\end{paytab}

由划线法可知唯一纯战略 NE 为 $(B,R)$。下面证明不存在非退化混合战略 NE。

设参与人 2 的混合战略为 $(q_L,q_C,q_R)$。参与人 1 各纯战略的期望收益为
\begin{align*}
  u_1(T)&=4q_C+5q_R,\\
  u_1(M)&=4q_L+5q_R,\\
  u_1(B)&=3q_L+3q_C+6q_R.
\end{align*}\indent
若 $T$ 与 $M$ 同以正概率出现，则需 $u_1(T)=u_1(M)$，即 $q_L=q_C$。此时
\[
  u_1(B)-u_1(T)=3q_L-q_C+q_R=2q_L+q_R.
\]\indent
若 $q_L+q_R>0$，则 $u_1(B)>u_1(T)$，$T$、$M$ 不是最佳反应；若 $q_L=q_R=0$（即 $q_C=1$），则 $u_1(T)=4>u_1(B)=3$，最佳反应只有 $T$。因此参与人 1 的任何混合 NE 支撑集不可能同时包含 $T$ 与 $M$。

同理，对称性表明参与人 2 的任何混合 NE 支撑集不可能同时包含 $L$ 与 $C$。

再检查其他可能的混合支撑：
\begin{enumerate}[leftmargin=2em]
  \item \textbf{参与人 1 以 $T,B$ 混合，参与人 2 以 $L,R$ 混合。} 参与人 1 对 $T$ 与 $B$ 无差异，要求
  \[
    4q_C+5q_R=3q_L+3q_C+6q_R
    \quad\Longrightarrow\quad
    q_C=3q_L+q_R.
  \]\indent
  参与人 2 对 $L$ 与 $R$ 无差异。令参与人 1 以概率 $p_T$ 选 $T$、以概率 $p_B=1-p_T$ 选 $B$，则
  \[
    u_2(L)=4p_T+3(1-p_T)=3+p_T,
    \qquad
    u_2(R)=3p_T+6(1-p_T)=6-3p_T.
  \]\indent
  令二者相等，得到
  \[
    p_T=\frac34,\qquad p_B=\frac14.
  \]\indent
  检查参与人 2 的第三个行动 $C$：
  \[
    u_2(C)=4p_M+3p_B=\frac34
    <u_2(L)=u_2(R)=\frac{15}{4}.
  \]\indent
  因而 $C$ 确实不在参与人 2 的支持集中，必须有 $q_C=0$。代回参与人 1 的无差异条件，
  \[
    0=q_C=3q_L+q_R.
  \]\indent
  因为概率非负，这迫使 $q_L=q_R=0$，与 $q_L+q_C+q_R=1$ 矛盾。因此这种支持组合不可能构成混合均衡。
  \item \textbf{参与人 1 以 $M,B$ 混合。} 由对称性，同理可排除。
  \item \textbf{参与人 1 只选 $T$ 或只选 $M$。} 若参与人 1 只选 $T$，参与人 2 的最佳反应为 $L$，但面对 $L$ 时参与人 1 会改选 $M$；若参与人 1 只选 $M$，参与人 2 的最佳反应为 $C$，但面对 $C$ 时参与人 1 会改选 $T$。两种情形都不能自洽。
\end{enumerate}

因此图 1.1.4 的唯一 NE 是纯战略 $(B,R)$，不存在混合战略 NE。

% ---------------------------------------------------------------------------
\exer{1.11}
解出习题 1.2 所给博弈的混合战略纳什均衡。

\sol
习题 1.2 的支付矩阵为
\begin{paytab}
\begin{tabular}{ccccc}
\multicolumn{2}{c}{} & \multicolumn{3}{c}{\pthead{参与人 2}}\\
\cline{3-5}
\multicolumn{2}{c}{} & $L$ & $C$ & $R$\\[2pt]
\hline
\multirow{3}{*}{参与人 1} & $T$ & $2,0$ & $1,1$ & $4,2$\\[4pt]
\cline{2-5}
 & $M$ & $3,4$ & $1,2$ & $2,3$\\[4pt]
\cline{2-5}
 & $B$ & $1,3$ & $0,2$ & $3,0$\\[4pt]
\hline
\end{tabular}
\end{paytab}

由习题 1.2 的分析，$B$ 与 $C$ 会被重复剔除严格劣战略剔除，因此任何混合 NE 中 $B$、$C$ 的概率均为 0，只需在支撑集 $\{T,M\}\times\{L,R\}$ 中寻找。

设参与人 1 以概率 $p$ 选 $T$、$1-p$ 选 $M$；参与人 2 以概率 $q$ 选 $L$、$1-q$ 选 $R$。

参与人 1 无差异条件：
\[
  2q+4(1-q)=3q+2(1-q)
  \quad\Longrightarrow\quad
  4-2q=2+q
  \quad\Longrightarrow\quad
  q=\frac23.
\]

参与人 2 无差异条件：
\[
  0\cdot p+4(1-p)=2p+3(1-p)
  \quad\Longrightarrow\quad
  4-4p=3-p
  \quad\Longrightarrow\quad
  p=\frac13.
\]

验证：$p=1/3\in[0,1]$，$q=2/3\in[0,1]$，满足概率约束。均衡中 $B$、$C$ 的概率为 0，也符合支撑集内的最佳反应。

因此混合战略纳什均衡为
\[
  \boxed{\sigma_1^*=\Bigl(\frac13,\frac23,0\Bigr),\qquad
          \sigma_2^*=\Bigl(\frac23,0,\frac13\Bigr)},
\]\indent
即参与人 1 以 $1/3$ 选 $T$、$2/3$ 选 $M$、不选 $B$；参与人 2 以 $2/3$ 选 $L$、$1/3$ 选 $R$、不选 $C$。

% ---------------------------------------------------------------------------
\exer{1.12}
求出下面标准式博弈的混合战略纳什均衡。

\begin{paytab}
\begin{tabular}{cccc}
\multicolumn{2}{c}{} & \multicolumn{2}{c}{\pthead{参与人 2}}\\
\cline{3-4}
\multicolumn{2}{c}{} & $L$ & $R$\\[2pt]
\hline
\multirow{2}{*}{参与人 1} & $T$ & $2,1$ & $0,2$\\[4pt]
\cline{2-4}
 & $B$ & $1,2$ & $3,0$\\[4pt]
\hline
\end{tabular}
\end{paytab}

\sol
设参与人 1 以概率 $p$ 选 $T$、$1-p$ 选 $B$；参与人 2 以概率 $q$ 选 $L$、$1-q$ 选 $R$。

\myfig{%
\resizebox{0.78\linewidth}{!}{%
\begin{tikzpicture}[every node/.style={font=\small}]
  \fill[lightbg] (0,1.05) rectangle (4.8,2.1);
  \draw[rulegray!80, line width=0.5pt] (0,0) rectangle (4.8,2.1);
  \draw[rulegray!80, line width=0.5pt] (2.4,0) -- (2.4,2.1);
  \draw[rulegray!80, line width=0.5pt] (0,1.05) -- (4.8,1.05);
  \node at (1.2,1.575) {$2,\,1$};
  \node at (3.6,1.575) {$0,\,2$};
  \node at (1.2,0.525) {$1,\,2$};
  \node at (3.6,0.525) {$3,\,0$};
  \node[left=7pt] at (0,1.575) {$T$};
  \node[left=7pt] at (0,0.525) {$B$};
  \node[above=6pt] at (1.2,2.1) {$L$};
  \node[above=6pt] at (3.6,2.1) {$R$};
  \node[above=19pt, font=\sffamily\bfseries] at (2.4,2.1) {参与人 2};
  \node[left=1.4cm, font=\sffamily\bfseries] at (0,1.05) {参与人 1};
  \node[right=9pt] at (4.8,1.575) {$p$};
  \node[right=9pt] at (4.8,0.525) {$1-p$};
  \node[below=7pt] at (1.2,0) {$q$};
  \node[below=7pt] at (3.6,0) {$1-q$};
\end{tikzpicture}}%
}{混合战略支付矩阵}

参与人 1 的无差异条件（$T$ 与 $B$ 期望收益相等）：
\[
  2q+0\cdot(1-q)=1\cdot q+3(1-q)
  \quad\Longrightarrow\quad
  2q=3-2q
  \quad\Longrightarrow\quad
  q=\frac34.
\]

参与人 2 的无差异条件（$L$ 与 $R$ 期望收益相等）：
\[
  1\cdot p+2(1-p)=2p+0\cdot(1-p)
  \quad\Longrightarrow\quad
  2-p=2p
  \quad\Longrightarrow\quad
  p=\frac23.
\]

验证：$p=2/3\in[0,1]$，$q=3/4\in[0,1]$，满足约束。因此混合战略纳什均衡为
\[
  \boxed{\sigma_1^*=\Bigl(\frac23,\frac13\Bigr),\qquad
          \sigma_2^*=\Bigl(\frac34,\frac14\Bigr)},
\]\indent
即参与人 1 以 $2/3$ 选 $T$、$1/3$ 选 $B$；参与人 2 以 $3/4$ 选 $L$、$1/4$ 选 $R$。

% ---------------------------------------------------------------------------
\exer{1.13}
两个企业各有一个工作空缺，假设企业所给的工资不同（其原因不在此处讨论，但关系到每一个空缺的价值）：企业 $i$ 给的工资为 $w_i$，这里 $(1/2)w_1<w_2<2w_1$。设想有两个工人，每人只能申请一份工作，两人同时决定是申请企业 1 的工作，还是向企业 2 申请。如果只有一个工人向一个企业申请，他就会得到这份工作；如果两个工人同时向一个企业申请工作，则企业随机选择一个工人，另一人就会失业（这时收益为 $0$）。解出两个工人标准式博弈的纳什均衡。（要更进一步了解企业是如何决定工资的，请参阅蒙哥马利（Montgomery），1991）

\sol
两工人的战略空间均为 $\{\text{企业 1},\text{企业 2}\}$。若两人同申企业 $i$，各得 $w_i/2$；若工人 1 申企业 1、工人 2 申企业 2，则收益为 $(w_1,w_2)$；反之则为 $(w_2,w_1)$。支付矩阵为

\begin{paytab}
\begin{tabular}{cccc}
\multicolumn{2}{c}{} & \multicolumn{2}{c}{\pthead{工人 2}}\\
\cline{3-4}
\multicolumn{2}{c}{} & 企业 1 & 企业 2\\[2pt]
\hline
\paybox{1.2cm}{工人 1} & \paybox{1.5cm}{企业 1} & \paybox{2.1cm}{$\dfrac{w_1}{2},\ \dfrac{w_1}{2}$} & \paybox{2.1cm}{$w_1,\ w_2$}\\
\hline
\paybox{1.2cm}{工人 1} & \paybox{1.5cm}{企业 2} & \paybox{2.1cm}{$w_2,\ w_1$} & \paybox{2.1cm}{$\dfrac{w_2}{2},\ \dfrac{w_2}{2}$}\\
\hline
\end{tabular}
\end{paytab}

\textbf{（1）纯战略纳什均衡。} 用划线法：
\begin{enumerate}[leftmargin=2em]
  \item 在（企业 1，企业 2）单元格，工人 1 得 $w_1$；若改申企业 2，得 $w_2/2<w_1$（因为 $w_2<2w_1$），不愿偏离。工人 2 得 $w_2$；若改申企业 1，得 $w_1/2<w_2$（因为 $w_1<2w_2$），也不愿偏离。故（企业 1，企业 2）是 NE。
  \item 同理，（企业 2，企业 1）也是 NE。
  \item 在（企业 1，企业 1）单元格，工人 1 得 $w_1/2$，改申企业 2 可得 $w_2>w_1/2$（因为 $2w_2>w_1$），会偏离，故不是 NE。同理（企业 2，企业 2）也不是 NE。
\end{enumerate}
因此纯战略 NE 为
\[
  \boxed{(\text{企业 1},\text{企业 2})\text{ 和 }(\text{企业 2},\text{企业 1})}.
\]

\textbf{（2）混合战略纳什均衡。} 设工人 1 以概率 $p$ 申企业 1、$1-p$ 申企业 2；工人 2 以概率 $q$ 申企业 1、$1-q$ 申企业 2。

工人 1 无差异条件：
\[
  q\cdot\frac{w_1}{2}+(1-q)w_1=q\,w_2+(1-q)\frac{w_2}{2}.
\]\indent
整理得
\[
  w_1-\frac{w_1}{2}q=\frac{w_2}{2}+\frac{w_2}{2}q
  \quad\Longrightarrow\quad
  q=\frac{2w_1-w_2}{w_1+w_2}.
\]

由对称性，工人 2 的无差异条件给出
\[
  p=\frac{2w_1-w_2}{w_1+w_2}.
\]

由 $(1/2)w_1<w_2<2w_1$ 可知 $0<p,q<1$，满足概率约束。因此混合战略 NE 为
\[
  \boxed{\sigma_1^*=\sigma_2^*=
    \Bigl(\frac{2w_1-w_2}{w_1+w_2},\ \frac{2w_2-w_1}{w_1+w_2}\Bigr)}.
\]\indent
每个工人以概率 $\dfrac{2w_1-w_2}{w_1+w_2}$ 申请企业 1，以概率 $\dfrac{2w_2-w_1}{w_1+w_2}$ 申请企业 2。

% ---------------------------------------------------------------------------
\exer{1.14}
证明附录 1.1.C 中的命题 B 不仅对纯战略成立，对混合战略同样成立：在混合战略纳什均衡中，概率大于 0 的战略一定不会被重复剔除严格劣战略所剔除。

\sol
设有限标准式博弈 $G=\{S_1,\dots,S_n;u_1,\dots,u_n\}$，$\sigma^*=(\sigma_1^*,\dots,\sigma_n^*)$ 是一个混合战略纳什均衡。对任意参与人 $i$，记 $S_i^+=\{s_i\in S_i:\sigma_i^*(s_i)>0\}$ 为 $i$ 在均衡中以正概率使用的纯战略集合。要证：每个 $s_i\in S_i^+$ 都 survive 重复剔除严格劣战略（IESDS）。

\textbf{引理。} 若 $s_i\in S_i^+$，则 $s_i$ 必为对手均衡混合战略 $\sigma_{-i}^*$ 的一个最佳反应（纯战略意义下）。

\textit{证明：} 由 NE 定义，$u_i(\sigma_i^*,\sigma_{-i}^*)\ge u_i(s_i',\sigma_{-i}^*)$ 对所有 $s_i'\in S_i$ 成立。而
\[
  u_i(\sigma_i^*,\sigma_{-i}^*)=\sum_{s_i\in S_i}\sigma_i^*(s_i)u_i(s_i,\sigma_{-i}^*).
\]\indent
若存在 $s_i\in S_i^+$ 使得 $u_i(s_i,\sigma_{-i}^*)<u_i(\sigma_i^*,\sigma_{-i}^*)$，则由于 $\sigma_i^*$ 以正概率加总所有 $s_i$，必有某个 $s_i'\in S_i^+$ 满足 $u_i(s_i',\sigma_{-i}^*)>u_i(\sigma_i^*,\sigma_{-i}^*)$，这与 NE 不等式矛盾。因此对所有 $s_i\in S_i^+$，$u_i(s_i,\sigma_{-i}^*)=u_i(\sigma_i^*,\sigma_{-i}^*)\ge u_i(s_i',\sigma_{-i}^*)$，即 $s_i$ 是最佳反应。$\square$

\textbf{主证明（数学归纳法）。} 对 IESDS 的轮次 $k=0,1,2,\dots$ 归纳证明：对每一 $i$，$\sigma_{-i}^*$ 仅在第 $k$ 轮后仍 surviving 的对手战略组合上赋予正概率。

\textit{基例 $k=0$：} 所有战略都在，结论显然成立。

\textit{归纳假设：} 设第 $k$ 轮后结论成立，即 $\sigma_{-i}^*$ 仅在第 $k$ 轮 surviving 的战略组合上赋予正概率。

\textit{归纳步骤：} 考虑第 $k+1$ 轮被剔除的某个战略 $s_i$。根据 IESDS 的定义，存在另一个（可能是混合的）战略 $\tilde\sigma_i$ 尚未被剔除，使得对所有由第 $k$ 轮 surviving 战略组成的对手战略组合 $s_{-i}$，都有
\[
  u_i(\tilde\sigma_i,s_{-i})>u_i(s_i,s_{-i}).
\]\indent
由期望的线性性，对任何仅在第 $k$ 轮 surviving 战略组合上赋予概率的混合战略 $\sigma_{-i}$，也有
\[
  u_i(\tilde\sigma_i,\sigma_{-i})>u_i(s_i,\sigma_{-i}).
\]\indent
特别地，取 $\sigma_{-i}=\sigma_{-i}^*$，由归纳假设它满足上述条件，故
\[
  u_i(\tilde\sigma_i,\sigma_{-i}^*)>u_i(s_i,\sigma_{-i}^*).
\]\indent
这说明 $s_i$ 不可能是 $\sigma_{-i}^*$ 的最佳反应。但引理告诉我们，每个 $s_i\in S_i^+$ 都是 $\sigma_{-i}^*$ 的最佳反应。因此 $s_i\notin S_i^+$，即 $s_i$ 在均衡中的概率为 0。

由于这对 IESDS 的每一轮都成立，故 $S_i^+$ 中的战略永远不会被剔除。命题 B 对混合战略纳什均衡同样成立。

\[
  \boxed{\text{在混合战略 NE 中，以正概率出现的纯战略必 survive 重复剔除严格劣战略。}}
\]\indent
